% Generated draft: requires mathematical and source review.
\documentclass[11pt]{article}
\usepackage[margin=25mm]{geometry}
\usepackage{fontspec}
\setmainfont{Latin Modern Roman}
\usepackage{amsmath,amssymb,mathtools,unicode-math}
\setmathfont{Latin Modern Math}
\usepackage{enumitem,needspace,xurl,hyperref}
\setlength{\parindent}{0pt}
\setlength{\parskip}{5pt}
\setlength{\emergencystretch}{3em}
\newcommand{\N}{\mathbb{N}}
\newcommand{\Z}{\mathbb{Z}}
\newcommand{\Q}{\mathbb{Q}}
\newcommand{\R}{\mathbb{R}}
\newcommand{\C}{\mathbb{C}}
\newcommand{\F}{\mathbb{F}}
\newcommand{\A}{\mathbb{A}}
\newcommand{\PP}{\mathbb{P}}
\newcommand{\E}{\mathbb{E}}
\newcommand{\eps}{\varepsilon}
\newcommand{\dd}{\,\mathrm{d}}
\begin{document}
\section*{UnsolvedMath ID: 30004335}
\textbf{Source problem code:} catalogue entry
\subsection{Definitions and mathematical statement}
Let \(K\) be a compact connected Lie group, and let \(\mathcal G_L(K)\) denote the set of all left-invariant Riemannian metrics on \(K\): for each \(g\in\mathcal G_L(K)\) and \(a\in K\), left translation \(L_a(x)=ax\) is an isometry. Write \(d_g\) for the Riemannian distance and \(\operatorname{vol}_g\) for the Riemannian volume measure. For \(x\in K\) and \(r>0\), define
\[
B_g(x,r)=\{y\in K:d_g(x,y)<r\},
\qquad
D(g)=\sup_{x\in K,\ r>0}
\frac{\operatorname{vol}_g(B_g(x,2r))}{\operatorname{vol}_g(B_g(x,r))}.
\]
The group \(K\) is called uniformly doubling if
\[
\sup_{g\in\mathcal G_L(K)}D(g)<\infty.
\]
\textbf{Question.} Is every compact connected Lie group uniformly doubling? Equivalently, is it true that for every compact connected Lie group \(K\) there exists a finite constant \(C_K\geq 1\) such that
\[
\operatorname{vol}_g(B_g(x,2r))
\leq C_K\operatorname{vol}_g(B_g(x,r))
\quad\text{for every }g\in\mathcal G_L(K),\ x\in K,\ r>0?
\]
The constant may depend on \(K\), but must be independent of the metric, center, and radius. All positive radii are included, with no restriction relative to the diameter. The supplied record labels this question open. Its background also mentions proved cases and unspecified ``related low-dimensional cases''; those are supplied literature qualifications, not additional hypotheses or independently verified claims.
\subsection{Short English statement}
For each compact connected Lie group \(K\), does there exist a finite bound, depending only on \(K\), for how much the volume of a ball can increase when its radius is doubled, simultaneously for every left-invariant Riemannian metric, every center, and every positive radius?
\subsection{Sources}
\begin{itemize}
\item Local UnsolvedMath registry, \texttt{data/unsolvedmath/problems.json}, permanent ID \(30004335\), source code \texttt{catalogue entry}, title ``Uniform Doubling for Compact Connected Lie Groups,'' in the supplied Oberwolfach Reports Open Problems collection. The record gives proposed year 2019 and names no individual proposer.
\item Supplied report citation: \emph{Heat Kernels, Stochastic Processes and Functional Inequalities} (2020), Oberwolfach Reports, pp.\ 3315--3407, DOI \texttt{10.4171/owr/2019/54}. \url{https://doi.org/10.4171/owr/2019/54}. The supplied background attributes the problem to workshop open-problem discussions documented in this report.
\item Supplied literature source: ``OWR report on heat kernels and functional inequalities.'' \url{https://ems.press/content/serial-article-files/46830?nt=1}.
\item Supplied literature source: ``Eldredge, uniform volume doubling for U(2).'' \url{https://www.birs.ca/workshops/2024/24w5314/files/Nathaniel%20Eldredge.pdf}.
\item Literature assessment dated exactly \texttt{2026-08-22}: the supplied curation labels the general conjecture open and reports uniform doubling proved for \(\mathrm{SU}(2)\), \(\mathrm{U}(2)\), and related low-dimensional cases. This is supplied dated curation derived from the OWR material, not an independently verified literature assessment; the record explicitly requests independent verification of time-sensitive status and cited claims.
\end{itemize}
\end{document}
